%%% =======================================================================
\section{Numerical Experiments}\label{sec:numerics}
%%% =======================================================================

% TODO: Replace with your numerical experiments. Typical contents:
%   - Experimental setup (domain, discretization, noise model, metrics)
%   - Test cases / phantoms
%   - Comparison with baseline methods
%   - Ablation studies
%   - Robustness / sensitivity analysis

%%% -----------------------------------------------------------------------
\subsection{Experimental Setup}\label{ssec:numerics_setup}
%%% -----------------------------------------------------------------------

Lorem ipsum dolor sit amet, consectetur adipiscing elit. Sed do eiusmod tempor
incididunt ut labore et dolore magna aliqua. Ut enim ad minim veniam, quis
nostrud exercitation ullamco laboris nisi ut aliquip ex ea commodo consequat.

\begin{table}[!htb]
    \centering
    \caption{Experimental parameters.}
    \label{tab:setup}
    \begin{tabular}{lc}
        \toprule
        Parameter & Value \\
        \midrule
        Domain size  & $N = 128$ \\
        Noise level  & $\delta = 1\%$ \\
        Regularization & $\beta = 10^{-3}$ \\
        \bottomrule
    \end{tabular}
\end{table}

%%% -----------------------------------------------------------------------
\subsection{Results}\label{ssec:results}
%%% -----------------------------------------------------------------------

Lorem ipsum dolor sit amet, consectetur adipiscing elit. Sed do eiusmod tempor
incididunt ut labore et dolore magna aliqua.

\begin{figure}[!htb]
    \centering
    % Example TikZ figure — replace with your own.
    % For external figures, use:  \includegraphics[width=0.75\textwidth]{your_figure.pdf}
    \begin{tikzpicture}[scale=1.2,
        box/.style={rectangle,draw,rounded corners,minimum width=2.2cm,minimum height=0.8cm,align=center},
        arr/.style={-{Stealth[scale=0.9]},thick}]
        % Forward problem
        \node[box,fill=blue!8] (model) at (0,0) {Unknown\\$u^\dagger$};
        \node[box,fill=green!8] (fwd) at (3.5,0) {Forward\\operator $\mathcal{A}$};
        \node[box,fill=red!8] (data) at (7,0) {Data\\$f = \mathcal{A}(u^\dagger)$};
        \draw[arr] (model) -- (fwd);
        \draw[arr] (fwd) -- (data);
        % Inverse problem
        \node[box,fill=red!8] (noisy) at (7,-2.2) {Noisy data\\$f^\delta = f + \eta$};
        \node[box,fill=green!8] (inv) at (3.5,-2.2) {Inverse\\solver};
        \node[box,fill=blue!8] (rec) at (0,-2.2) {Reconstruction\\$u_h \approx u^\dagger$};
        \draw[arr] (noisy) -- (inv);
        \draw[arr] (inv) -- (rec);
        % Noise arrow
        \draw[arr,dashed,gray] (data) -- (noisy) node[midway,right,gray] {$\eta$};
    \end{tikzpicture}
    \caption{Schematic of the forward and inverse problem. The forward operator
             $\mathcal{A}$ maps the unknown $u^\dagger$ to noise-free data $f$;
             the inverse solver recovers an approximation $u_h$ from noisy
             measurements $f^\delta = f + \eta$.}
    \label{fig:results}
\end{figure}

\begin{table}[!htb]
    \centering
    \caption{Quantitative comparison. Mean $\pm$ std over five realizations.}
    \label{tab:results}
    \begin{tabular}{lccc}
        \toprule
        Method   & Metric 1 & Metric 2 & Metric 3 \\
        \midrule
        Baseline & -- & -- & -- \\
        Proposed & -- & -- & -- \\
        \bottomrule
    \end{tabular}
\end{table}

%%% -----------------------------------------------------------------------
\subsection{Discussion}\label{ssec:discussion}
%%% -----------------------------------------------------------------------

Lorem ipsum dolor sit amet, consectetur adipiscing elit. Sed do eiusmod tempor
incididunt ut labore et dolore magna aliqua. Duis aute irure dolor in
reprehenderit in voluptate velit esse cillum dolore eu fugiat nulla pariatur.
